\section{Preliminaries and Analysis}
\label{sec:analysis}

% ---- 脉络（写完删除）--------------------------------------------------------
% 目的：先用实验把 video model 的"局限"变成事实，使 Method 成为顺理成章的回应。
% 1. 设定：WAM 结构（observation/future token、读出层、flow-matching action decoder）、
%    token 账（2,400 obs + 16,800 future）、部署运行点（单次截断前向，tau_v = 1）。
% 2. 冗余有多少：预算曲线。
% 3. 为什么相似度失效：部署点上相邻帧相似度不可复现；随机排序与相似度排序等效；
%    future 帧的 "motion" 随噪声变化而非随观测变化。
% 4. 为什么删除失效：删除式剪枝在激进预算下崩溃，合并平稳。
% 5. 位置约定在两个轴上相反：空间保留原坐标，时间必须紧凑。
%    （"时间位置是槽位而非内容"是解释，写成"一致于"，不写"证明"。）
% 6. 小结 -> 三条设计原则，引出 Sec.~\ref{sec:method}。
% ---------------------------------------------------------------------------

\TODO{Preliminaries and analysis.}
